\section{Claim Scope and Paired Comparisons}
\label{app:claimscope}

The deterministic topological derivative and the off-diagonal product
construction are established ingredients. The claimed contribution is their
use as a direct grid-free insertion query, the accompanying error accounting,
and its evaluated integration into finite geometric decisions. The exact
sphere check in \cref{app:exactcavity} is verification evidence, not a claim
that the radial Poisson solution is new.

\paragraph{Three different levels of guarantee.}
The direct-score unbiasedness and variance statements concern ideal marginal
PDE oracles and independent complete pairs. Finite-tolerance runs approximate
those oracles. The general-objective insertion theorem additionally requires
numerically usable bounds on marginal bias and the finite-radius remainder;
we have not supplied those bounds for every displayed mesh and objective.
The specialized Poisson-compliance validator instead avoids an unknown
remainder through a variational bound and handles truncation through a
one-sided construction. Its probability guarantee still assumes the stated
conditional sampling laws; the floating-point implementation is not a
verified sampler. The design pipelines and channel-placement heuristic are
not certified by either theorem. Confidence bands in the plots quantify
sampling uncertainty, not those missing numerical error budgets.

\paragraph{Paired differences distinguish placement value from solver value.}
\Cref{tab:pairedclaimscope} uses the existing paired bootstrap outputs,
without selecting new seeds or rerunning optimization. Differences compare
TWoS after eight refinement rounds with multigrid/proximity after eight,
and hottest-point placement after sixteen. Positive values favor TWoS.
The intervals are pointwise and exclude numerical bias; they are not
simultaneous comparisons or equivalence tests. An interval containing zero
does not prove that the methods are equivalent. The uncertainty and reference
limitations in the main text apply to every entry.

\begin{table}[t]
\centering\small
\caption{Paired mean cooling difference, in percentage points, with pointwise
95\% joint bootstrap intervals. MG uses longest-axis grid count 65.}
\label{tab:pairedclaimscope}
\begin{tabular}{lrrr}
\toprule
Geometry & TWoS minus MG & TWoS minus nearest & TWoS minus hottest\\
\midrule
Finned housing & -0.11 [-1.24, +0.99] & +1.43 [-0.32, +3.39] & +22.25 [+12.76, +37.94] \\
Branching manifold & +0.14 [-0.74, +1.63] & -0.93 [-1.40, -0.55] & +16.26 [+15.22, +17.43] \\
Annular spreader & -0.32 [-1.11, +0.18] & -0.84 [-1.99, +0.66] & +11.35 [+10.39, +12.24] \\
\bottomrule
\end{tabular}
\end{table}
